\chapter{Các số tập trung đo mức độ phân tán của mẫu số liệu ghép nhóm}

\section{Khoảng biến thiên và khoảng tứ phân vị}

\exercise Khảo sát thời gian tự học (phút/ngày) của 40 học sinh lớp 12, có được bảng số liệu ghép nhóm:
\begin{center}
\begin{tblr}{
  hlines, vlines,
  colspec = {ccccc},
}
Thời gian (phút) & $[30; 60)$ & $[60; 90)$ & $[90; 120)$ & $[120; 150)$ \\
Số học sinh & $8$ & $15$ & $12$ & $5$ \\
\end{tblr}
\end{center}
Hãy xác định khoảng biến thiên $R$ của mẫu số liệu ghép nhóm trên.

\exercise Khảo sát cân nặng (kg) của 50 học sinh trong một lớp, thu được bảng số liệu ghép nhóm:
\begin{center}
\begin{tblr}{
  hlines, vlines,
  colspec = {cccccc},
}
Cân nặng (kg) & $[40; 45)$ & $[45; 50)$ & $[50; 55)$ & $[55; 60)$ & $[60; 65)$ \\
Số học sinh & $6$ & $14$ & $18$ & $8$ & $4$ \\
\end{tblr}
\end{center}
Xác định nhóm chứa tứ phân vị thứ nhất $Q_1$, nhóm chứa tứ phân vị thứ ba $Q_3$, và tính khoảng tứ phân vị $\Delta_Q$ của mẫu số liệu ghép nhóm đó.

\exercise Bảng thống kê tuổi thọ (giờ) của 100 bóng đèn LED do một công ty sản xuất thử nghiệm như sau:
\begin{center}
\begin{tblr}{
  hlines, vlines,
  colspec = {cccccc},
}
Tuổi thọ (giờ) & $[1000; 1200)$ & $[1200; 1400)$ & $[1400; 1600)$ & $[1600; 1800)$ & $[1800; 2000)$ \\
Số bóng đèn & $10$ & $25$ & $40$ & $20$ & $5$ \\
\end{tblr}
\end{center}
Tính khoảng tứ phân vị $\Delta_Q$ của mẫu số liệu ghép nhóm trên và giải thích ý nghĩa của giá trị thu được đối với độ phân tán của $50\%$ số liệu chính giữa.

\exercise Khảo sát thời gian truy cập Internet mỗi ngày (giờ) của học sinh ở hai trường A và B cho kết quả:
\begin{center}
\begin{tblr}{
  hlines, vlines,
  colspec = {ccccc},
}
Thời gian (giờ) & $[0; 1)$ & $[1; 2)$ & $[2; 3)$ & $[3; 4)$ \\
Số HS Trường A & $15$ & $35$ & $30$ & $20$ \\
Số HS Trường B & $5$ & $20$ & $55$ & $20$ \\
\end{tblr}
\end{center}
Hãy tính khoảng tứ phân vị $\Delta_Q$ cho mẫu số liệu ghép nhóm của từng trường. Dựa trên kết quả đó, hãy so sánh độ phân tán thời gian truy cập Internet ở $50\%$ số liệu chính giữa của học sinh hai trường và đưa ra nhận xét.

\section{Phương sai và độ lệch chuẩn}

\exercise Cho mẫu số liệu ghép nhóm về chiều cao (cm) của 30 cây keo giống:
\begin{center}
\begin{tblr}{
  hlines, vlines,
  colspec = {ccccc},
}
Chiều cao (cm) & $[20; 30)$ & $[30; 40)$ & $[40; 50)$ & $[50; 60)$ \\
Số cây & $5$ & $12$ & $10$ & $3$ \\
\end{tblr}
\end{center}
Xác định giá trị đại diện $x_i$ của mỗi nhóm và tính số trung bình $\bar{x}$ của mẫu số liệu ghép nhóm.

\exercise Thống kê số tiền điện (nghìn đồng) phải trả trong tháng 8 của 40 hộ gia đình trong một khu phố được cho bởi bảng sau:
\begin{center}
\begin{tblr}{
  hlines, vlines,
  colspec = {ccccc},
}
Số tiền (nghìn đồng) & $[200; 400)$ & $[400; 600)$ & $[600; 800)$ & $[800; 1000)$ \\
Số hộ gia đình & $6$ & $14$ & $12$ & $8$ \\
\end{tblr}
\end{center}
Tính phương sai $s^2$ và độ lệch chuẩn $s$ của mẫu số liệu ghép nhóm trên.

\exercise Kết quả kiểm tra môn Toán của hai lớp 12A và 12B (mỗi lớp có 40 học sinh) được thu thập ở bảng dưới đây:
\begin{center}
\begin{tblr}{
  hlines, vlines,
  colspec = {ccccc},
}
Điểm số & $[4; \num{5.5})$ & $[\num{5.5}; 7)$ & $[7; \num{8.5})$ & $[\num{8.5}; 10]$ \\
Số HS Lớp 12A & $4$ & $12$ & $18$ & $6$ \\
Số HS Lớp 12B & $8$ & $8$ & $10$ & $14$ \\
\end{tblr}
\end{center}
\begin{enumerate}
    \item Tính điểm trung bình và độ lệch chuẩn về điểm số của từng lớp.
    \item Dựa vào độ lệch chuẩn, cho biết học sinh lớp nào có điểm thi đồng đều hơn.
\end{enumerate}

\exercise Một nhà máy sản xuất linh kiện điện tử kiểm tra bán kính (mm) của các sản phẩm xuất xưởng từ hai dây chuyền sản xuất I và II. Bảng số liệu ghép nhóm thu được như sau:
\begin{center}
\begin{tblr}{
  hlines, vlines,
  colspec = {ccccc},
}
Bán kính (mm) & $[\num{9.8}; \num{9.9})$ & $[\num{9.9}; \num{10.0})$ & $[\num{10.0}; \num{10.1})$ & $[\num{10.1}; \num{10.2})$ \\
Dây chuyền I & $10$ & $40$ & $40$ & $10$ \\
Dây chuyền II & $25$ & $25$ & $25$ & $25$ \\
\end{tblr}
\end{center}
\begin{enumerate}
   \item Tính số trung bình và độ lệch chuẩn của bán kính sản phẩm ở từng dây chuyền.
   \item Dựa vào các số đặc trưng đo mức độ trung tâm và độ phân tán, hãy đánh giá dây chuyền nào vận hành ổn định hơn và ít nguy cơ tạo ra sản phẩm lỗi hơn.
\end{enumerate}